%! AP Statistics — Unit 1, Lesson 1.10 — Group Activity
\documentclass[10pt]{article}
\usepackage{apstats-article}
\usepackage{apstats-boxes}

%=============================================================================
\begin{document}

\pageheader{Unit 1, Lesson 1.10}{Group Activity: The Normal Distribution}
\namepartnerperiod

\vspace{0.10in}

\begin{tcolorbox}[colback=sky, colframe=navy, boxrule=0.8pt, arc=2mm,
    left=3mm, right=3mm, top=2mm, bottom=2mm]
\small
\textbf{Part 1} (8 min): Apply the 68--95--99.7 Rule.\\
\textbf{Part 2} (8 min): Compute $z$-scores and use the Normal table.\\
\textbf{Part 3} (6 min): Find a percentile and assess an unusual value.
\end{tcolorbox}

\vspace{0.10in}

% ── CONTEXT ──────────────────────────────────────────────────────────────────
\begin{tcolorbox}[colback=goldbg, colframe=goldacc, boxrule=0.7pt, arc=2mm,
    left=3mm, right=3mm, top=2mm, bottom=2mm]
\small
\textbf{Context:} SAT scores for a large population of college-bound students are
approximately Normally distributed with mean $\mu = 1060$ and standard deviation
$\sigma = 195$. Write this distribution using notation: \blank{4cm}
\end{tcolorbox}

\vspace{0.12in}

% ── PART 1 ───────────────────────────────────────────────────────────────────
{\large\bfseries\color{navy} Part 1 \quad Applying the 68--95--99.7 Rule}

\vspace{0.10in}

\begin{enumerate}[leftmargin=1.5em, itemsep=8pt]

  \item Fill in the boundaries for each Empirical Rule region.\\[6pt]
        \renewcommand{\arraystretch}{1.5}
        \begin{tabularx}{\linewidth}{@{} l Y Y Y @{}}
        \rowcolor{sky}
        Region & Lower bound & Upper bound & \% of scores \\
        \midrule
        $\mu \pm 1\sigma$ & \blank{2cm} & \blank{2cm} & \blank{2cm}\% \\
        $\mu \pm 2\sigma$ & \blank{2cm} & \blank{2cm} & \blank{2cm}\% \\
        $\mu \pm 3\sigma$ & \blank{2cm} & \blank{2cm} & \blank{2cm}\% \\
        \end{tabularx}

  \item What percentage of students scored \textbf{above} 1255?
        Show your reasoning using the Empirical Rule.\\[5pt]
        \writeline\\[1pt]\writeline

  \item What percentage of students scored \textbf{below} 670?\\[5pt]
        \writeline\\[1pt]\writeline

  \item What SAT score is at the \textbf{84th percentile}?
        (\textit{Hint}: 84\% $=$ 50\% below the mean $+$ half of the 68\% region.)\\[5pt]
        \writeline

\end{enumerate}

\vspace{0.12in}

% ── PART 2 ───────────────────────────────────────────────────────────────────
{\large\bfseries\color{navy} Part 2 \quad $z$-Scores and the Normal Table}

\vspace{0.10in}

\begin{enumerate}[leftmargin=1.5em, itemsep=8pt, resume]

  \item Calculate the $z$-score for a student who scored \textbf{900}.
        Show work and interpret in context.\\[5pt]
        $z = \dfrac{x - \mu}{\sigma} = \dfrac{\blank{1.5cm} - \blank{1.5cm}}{\blank{1.5cm}} = \blank{2cm}$\\[6pt]
        Interpretation: \writeline

  \item Use the Normal table (Table A) to find $P(\text{SAT} < 900)$.\\[5pt]
        $P(Z < \blank{2cm}) = \blank{2cm}$\\[4pt]
        In context: approximately \blank{3cm}\% of students scored below 900.

  \item Find $P(\text{SAT} > 1400)$.
        First compute the $z$-score, then use Table A.\\[5pt]
        $z = \dfrac{1400 - 1060}{195} = \blank{2cm}$\\[4pt]
        $P(Z > \blank{2cm}) = 1 - P(Z < \blank{2cm}) = 1 - \blank{2cm} = \blank{2cm}$

  \item Find $P(900 < \text{SAT} < 1400)$.\\[5pt]
        $P(900 < \text{SAT} < 1400) = P(Z < \blank{2cm}) - P(Z < \blank{2cm}) = \blank{2cm} - \blank{2cm} = \blank{2cm}$

\end{enumerate}

\vspace{0.12in}

% ── PART 3 ───────────────────────────────────────────────────────────────────
{\large\bfseries\color{navy} Part 3 \quad Percentiles and Unusual Values}

\vspace{0.10in}

\begin{enumerate}[leftmargin=1.5em, itemsep=8pt, resume]

  \item Find the SAT score at the \textbf{90th percentile} using the inverse Normal.\\[5pt]
        Step 1: Find $z^*$ such that $P(Z < z^*) = 0.90$. From Table A: $z^* \approx \blank{2cm}$\\[4pt]
        Step 2: $x = \mu + z^* \cdot \sigma = 1060 + \blank{2cm} \times 195 = \blank{3cm}$\\[4pt]
        Interpretation: A student would need approximately \blank{2cm} to be in the top 10\%.

  \item A student claims to have scored \textbf{1500} on the SAT.
        Is this score unusually high? Justify with a $z$-score.\\[5pt]
        $z = \dfrac{1500 - 1060}{195} = \blank{2cm}$\\[5pt]
        \writeline\\[1pt]\writeline

  \item \textbf{Discuss with your group:} Would it be appropriate to use the Normal model
        for a class of 20 students if you saw that the histogram was right-skewed with one
        very high outlier? Explain.\\[5pt]
        \writeline\\[1pt]\writeline

\end{enumerate}

\end{document}
